\section{Conclusion}
\label{sec:conclusion}

Enterprise LLM agents treat every retrieved item as both sendable and
trustworthy the moment it enters the context window. Neither property is
conferred by retrieval, and the second one has had a formalism since 1977 that
the agent-privacy literature has not used.

We presented \sys, a zero-trust policy layer at the agent's write boundary
that labels each retrieved entry on two independent axes and enforces
Bell--LaPadula no-write-down together with Biba no-write-up per request.
Three findings stand out. The axes are genuinely independent, and the
essential/sensitive partition that existing benchmarks rely on is not merely
incomplete but unsound, since it cannot represent an entry that is at once
required by the task and forbidden to the recipient. Integrity for an agent is
a capability constraint rather than a filtering problem: bounding the
integrity of the action by the integrity of the weakest supporting evidence
requires no natural-language inference and cannot be overridden by the
context. And enforcement placement dominates enforcement strength, with the
same policy run before generation recovering all of the utility that post-hoc
filtering loses.

The direction we think matters most is the one our own metrics cannot yet
reach. Integrity-Leakage measures whether corrupting context was
\emph{repeated}, and we showed by measurement that this undercounts the failure
by exactly the entries that were silently \emph{believed}. Measuring reliance
rather than disclosure---through paired counterfactuals and per-claim
grounding---is what would turn the integrity axis from a proxy into a
measurement, and it is where we are taking this work next.

\section{Ethics and Open Science}
\label{sec:ethics}

\subsection{Ethical considerations}

\paragraph{Data.} All benchmark cases are synthetic. No real customer,
employee, or organizational data appears in any seed. Entries that model
regulated content---named-customer incidents, personnel matters---are
fabricated for the purpose and are marked \texttt{regulated} in the fixture so
that the hard-block path is exercised. We inherit this posture from CI-Work,
which uses synthetic enterprise data for the same reason.

\paragraph{Dual use.} The benchmark contains a class of entries
($\Ecorr$) designed to corrupt an agent's verdict, which is a mild
dual-use concern: the construction recipe is also a recipe for adversarial
context. We judge the disclosure worthwhile on the standard grounds that the
technique is already documented in the indirect prompt injection
literature~\cite{greshake2023indirect}, the entries are deliberately
low-sophistication, and a defense cannot be evaluated against attacks that are
not published. Our corrupting entries assert unverifiable claims in channels
the adversary legitimately controls; they contain no novel attack technique.

\paragraph{Incident reporting.} \Cref{tab:incidents} discusses publicly
reported incidents at named organizations. We rely exclusively on published
accounts, several of which are the affected parties' own disclosures, and we
draw no conclusions about culpability. Our interest is the structural pattern.

\paragraph{Deployment risk of the work itself.} A policy layer that reports
zero leakage on a benchmark may induce misplaced confidence in deployment. We
have tried to mitigate this in the writing: every headline number is reported
with the condition under which it fails (\Cref{sec:eval:rq2}), the integrity
metric's blind spot is reported as a measured negative result
(\Cref{sec:eval:ilblindspot}), and the inline arm is labeled a policy ceiling
rather than an achieved score wherever it appears.

\subsection{Open science}
\label{sec:ethics:openscience}

We will release the policy layer, the benchmark fixtures, the configurations,
and the evaluation harness under an open license. Three properties are
deliberate and we consider them part of the contribution:

\begin{enumerate}
  \item \textbf{The offline path has no dependencies and no network access.}
  The unit tests, both evaluation runs, and the verifier execute in under a
  second, so a reviewer can reproduce every offline number without credentials
  or cost.
  \item \textbf{Every reported number traces to a configuration file.} No
  weight is hard-coded in the implementation. The parameter tiering of
  \Cref{sec:method:weights} is expressed in the configuration, so a reader can
  see which values are structural, which are calibrated, and which are
  operating-point choices.
  \item \textbf{A verifier re-derives the paper's decisions from the
  configuration.} It recomputes the policy decisions for all three request
  configurations reported in \Cref{sec:eval} and asserts them, failing if the
  prose and the implementation diverge. It additionally asserts the safety
  property of \Cref{sec:discussion:bugs}, that no regulated entry is released
  by the threshold override at any operating point.
\end{enumerate}

The LLM-judge path requires API credentials and is therefore not reproducible
without cost; we will report the judge model and version, and release the
prompts and the raw judge outputs so that the scoring is auditable even where
it is not re-runnable.

\paragraph{Generative AI disclosure.} \todo{State the extent of generative AI
use in preparing this work, per venue policy. SaTML/ACM both require an
explicit statement even when no tools were used.}
